% !TEX root = ../Thesis.tex

\chapter{Introduzione}
L'esigenza intrinseca e crescente di garantire comunicazioni che siano contemporaneamente riservate, autentiche e integre rappresenta un pilastro fondamentale su cui si erge l'architettura stessa del mondo digitale contemporaneo. La sicurezza delle informazioni, in passato, è stata prevalentemente affidata a robuste tecniche crittografiche e a protocolli standard ben consolidati, quali il Transport Layer Security (TLS) e il suo predecessore Secure Sockets Layer (SSL). Questi sistemi sono stati meticolosamente progettati per rendere il contenuto effettivo dei messaggi inaccessibile a qualsiasi entità esterna non in possesso delle chiavi di decifratura appropriate; nel contesto di TLS, ciò si traduce tipicamente nell'esclusione di chiunque non partecipi al processo di negoziazione della connessione (l'handshake).

Tuttavia, l'escalation senza precedenti delle misure di sorveglianza di massa su scala globale e la proliferazione di meccanismi di censura sempre più capillari e sofisticati hanno messo a nudo una vulnerabilità critica, intrinseca e finora difficile da eradicare alla crittografia moderna: la sua stessa visibilità. Sebbene la crittografia riesca efficacemente a mantenere segreto il contenuto dei dati trasmessi, l'elevata entropia e le caratteristiche distintive del traffico cifrato lo rendono facilmente identificabile e distinguibile dalle comunicazioni in chiaro. Questo stato di cose consente a regimi autoritari, a entità statali o a sofisticati sistemi di filtraggio di rete (come i firewall avanzati) di individuare, bloccare o degradare selettivamente le comunicazioni protette, compromettendo così l'effettivo scambio di informazioni sensibili in contesti operativi ostili o sorvegliati.

Di fronte a tali sfide, specialmente in scenari caratterizzati da censura estrema o da ambienti di comunicazione altamente avversari, emerge in modo prepotente la necessità di esplorare e impiegare metodologie steganografiche. La steganografia, come disciplina, si dedica all'arte e alla scienza dell'occultamento di informazioni segrete all'interno di oggetti di copertura (cover objects) che, a un'osservazione superficiale, appaiono del tutto innocui e non sospetti. Questi oggetti di copertura possono assumere svariate forme, tra cui testi, immagini, file audio o video. A differenza della crittografia, il cui successo e la cui sicurezza dipendono in larga misura dalla complessità matematica o computazionale necessaria per decifrare il messaggio cifrato, la sicurezza in ambito steganografico poggia su un principio radicalmente diverso: l'indistinguibilità statistica. L'obiettivo primario è fare in modo che un osservatore esterno non sia minimamente in grado di discriminare tra un oggetto che contiene un messaggio segreto (denominato stego-object) e un oggetto analogo che ne è privo.

Le metodologie steganografiche classiche, quali ad esempio la manipolazione dei bit meno significativi (Least Significant Bit - LSB) nelle immagini digitali, si sono progressivamente rivelate vulnerabili agli attacchi di crittoanalisi statistica sempre più avanzati e raffinati. Per superare queste limitazioni intrinseche e garantire un livello di sicurezza che sia non solo robusto ma anche dimostrabile, la ricerca più recente ha abbracciato un vero e proprio cambio di paradigma. Questo nuovo orientamento si concentra sull'utilizzo strategico di modelli generativi, in particolare quelli basati su reti neurali profonde.

In questo fertile contesto di ricerca si inserisce \textbf{Meteor}, un innovativo algoritmo di steganografia a chiave simmetrica. Meteor sfrutta le potenti capacità computazionali e generative dei modelli di linguaggio di grandi dimensioni (Large Language Models - LLM), originariamente derivati da architetture come GPT-2, per creare canali di comunicazione occulti e resilienti. A differenza delle tecniche steganografiche tradizionali che modificano un contenuto esistente, Meteor interviene e guida il processo di generazione stessa del contenuto. Esso utilizza la distribuzione di probabilità intrinseca del modello generativo per "pilotare" il campionamento dei token successivi in base al messaggio segreto opportunamente cifrato. Questo approccio garantisce che lo stego-object risultante sia statisticamente coerente e indistinguibile dall'output naturale del modello, rendendo la presenza del messaggio nascosto estremamente difficile, se non impossibile, da rilevare attraverso analisi statistiche convenzionali.

L'ambizioso obiettivo della presente tesi triennale è estendere il principio fondamentale di sicurezza steganografica ideato da Meteor, originariamente applicato al dominio testuale, all'ambito delle immagini digitali ad alta risoluzione.
Per conseguire questo scopo, viene sfruttata l'architettura \textbf{VQ-GAN} (Vector Quantized Generative Adversarial Network). La scelta strategica di questa specifica architettura deriva dalla sua intrinseca capacità di modellare immagini estremamente complesse attraverso la creazione e l'utilizzo di un vocabolario discreto di "token visivi", memorizzati in un codebook. Questa rappresentazione discreta apre le porte all'applicazione di un modello Transformer autoregressivo, ampiamente utilizzato nella generazione sequenziale, per guidare il processo di sintesi dell'immagine.
Il lavoro svolto mira dunque a dimostrare con rigore e sperimentazione come sia possibile incorporare messaggi cifrati direttamente durante la fase di generazione dell'immagine da parte dell'architettura VQ-GAN. L'obiettivo finale è ottenere immagini non solo visivamente simili a quelle generate dalla stesso modello, ma anche intrinsecamente robuste contro tentativi sofisticati di rilevamento steganografico.

Le implicazioni pratiche e teoriche di questo approccio sono potenzialmente rivoluzionarie per lo sviluppo di strumenti di comunicazione \emph{censorship-resistant}. La capacità di abilitare lo scambio sicuro e nascosto di informazioni all'interno di immagini di alta qualità offre un livello inedito di protezione per la privacy degli utenti e per la libertà di espressione in contesti digitali sempre più soggetti a sorveglianza e controllo.

\bigskip

Il presente elaborato è stato organizzato con cura per condurre il lettore attraverso un percorso logico e approfondito, dalla teoria alla pratica:

\begin{itemize}
	\item Il \textbf{\autoref{ch:stato-arte}} fornisce un'analisi completa e dettagliata del background teorico indispensabile, approfondendo lo stato dell'arte delle tecniche steganografiche esistenti e i principi fondamentali di funzionamento dei modelli generativi più avanzati, con un'enfasi particolare sulle architetture GAN (Generative Adversarial Networks) e Transformer.
	\item Il \textbf{\hyperref[cap-3]{Capitolo 3}} presenta le definizioni formali e le specifiche tecniche del sistema proposto, illustrando in dettaglio l'architettura implementata e le modifiche apportate al meccanismo di campionamento dei token visivi per consentire l'embedding del messaggio segreto cifrato direttamente nel flusso generativo. Vengono inoltre discussi i limiti teorici e pratici del lavoro svolto, evidenziando le sfide affrontate e le soluzioni adottate per superarle.
	
	\item Il \textbf{\hyperref[cap-4]{Capitolo 4}} presenta l'idea alla base del lavoro di tesi,l'architettura completa del sistema implementato e i limiti teorici e pratici del lavoro svolto. Vengono illustrate in dettaglio le modifiche apportate al meccanismo di campionamento dei token visivi per consentire l'embedding del messaggio segreto cifrato direttamente nel flusso generativo.
	
	\item Il \textbf{\hyperref[cap-5]{Capitolo 5}} espone in modo esaustivo i risultati delle sperimentazioni condotte. Vengono valutate sia la qualità visiva delle immagini generate, attraverso l'utilizzo di metriche consolidate come la Fréchet Inception Distance (FID), sia la capacità steganografica del sistema, analizzando il payload trasferibile e la robustezza contro attacchi di rilevamento.
	
	\item Infine, il \textbf{\hyperref[cap-6]{Capitolo 6}} trae le conclusioni finali del lavoro svolto, riassumendo i contributi principali della ricerca e delineando una serie di possibili sviluppi futuri, mirati al miglioramento incrementale dell'efficienza, della robustezza e della sicurezza complessiva del sistema proposto.
\end{itemize}